\honest{Optimization and the Intelligence Explosion}{Eliezer Yudkowsky}{2015}

A mind's power is its ability to hit small targets in a large space of possibilities. A car is improbable among random arrangements of its atoms; natural selection and human intelligence both produce such improbable things, so I call both optimization processes. I grant at once that this is an approximation.

I separate the meta level, what does the optimizing, from the object level, what gets optimized. Natural selection is simple and began by accident; a cat is far more complicated. Animal brains learned, but their learning did not accumulate, so they were pieces in the game, not players, until human brains very recently became players.

The meta level has almost always been protected. Genes rarely change natural selection; the exceptions, sex and cells and DNA, are rare enough that ``evolutionary biologists structure entire histories of life on Earth around them.'' \nb{The best-known such history is Maynard Smith and Szathm\'ary's \textsc{The Major Transitions in Evolution} (1995), which includes the origin of chromosomes, of DNA with proteins as enzymes, of sex, and of language. This supports the post.} Science builds on science, but the brain doing it is the one the first farmers had. That rationality training cannot turn an arbitrary person into Einstein ``shows the power of a few minor genetic quirks of brain design.'' \nb{It shows that training does not close the gap. That the gap consists of genetic quirks is assumed.}

Now, some of us want to build an intelligence that can redesign itself all the way down. It would have no protected level that optimizes, and that would break with the whole past. \nb{The idea of a machine that designs better machines, and so sets off an ``intelligence explosion,'' is I.~J. Good's, from 1965. The title uses Good's term; the post does not mention Good.}

Until now optimizers have worked ``at a constant rate,'' and history sped up because each innovation opened the way to others, and because, once in a long while, something like sex or science reached the meta level. \nb{The number of optimizers grew. Robin Hanson wrote in 2008 that growth also speeds up ``as the economy, by getting larger, enables its members to explore an ever-increasing number of innovations.''} People ask what happens if each linear gain needs exponentially more self-rewriting. My answer: natural selection exerted ``roughly constant optimization power'' on the hominid line, and ``this doesn't seem to have required exponentially more time for each linear increment of improvement.'' \nb{Neither the rate nor the increment is measured, and the case covers only the climb from ape to human. The objection concerns an AI improving itself far beyond that.}

I admit that all this is ``mere analogic reasoning,'' good ``at best'' for qualitative predictions. Still, it is my reason for not extending the graph of growth over time past the arrival of such an AI: the right graph plots optimization power in against optimized product out. \nb{The first version of this post, from June 2008, named the forecast it answers: Robin Hanson's, that a new transition could make the world economy double in a week to a month.}
